\documentclass[12pt,a4paper]{article}
\usepackage[utf8]{inputenc}
\usepackage[T1]{fontenc}
\usepackage[french]{babel}
\usepackage{graphicx}
\usepackage{hyperref}
\usepackage{geometry}
\usepackage{booktabs}
\usepackage{float}
\usepackage{amsmath}

\geometry{
    a4paper,
    total={170mm,257mm},
    left=20mm,
    top=20mm,
}

\title{Rapport Final : Apprentissage Fédéré pour la Détection des Maladies Cardiaques}
\author{Équipe de Projet}
\date{\today}

\begin{document}

\maketitle

\section{Introduction}
Ce rapport présente les résultats du projet d'Apprentissage Fédéré (Federated Learning - FL) appliqué à la détection des maladies cardiaques. L'objectif était de simuler un environnement où 8 hôpitaux collaborent pour entraîner un modèle de Machine Learning (Multi-Layer Perceptron) sans partager leurs données médicales locales.

L'étude compare quatre optimisateurs répartis :
\begin{itemize}
    \item \textbf{Local SGD} : Entraînement local sans agrégation (Baseline).
    \item \textbf{FedAvg} : Moyennage fédéré standard.
    \item \textbf{FedAdam} : Optimisation côté serveur utilisant l'algorithme Adam.
    \item \textbf{FedProx} : Ajout d'un terme proximal pour gérer l'hétérogénéité des données locales.
\end{itemize}

Les expériences ont été menées sur deux partitions de données :
\begin{itemize}
    \item \textbf{IID} (Indépendantes et Identiquement Distribuées) : Chaque hôpital a une distribution équilibrée des classes.
    \item \textbf{Non-IID} (Dirichlet avec $\alpha=0.5$) : Déséquilibre fort simulant des hôpitaux spécialisés ou géographiquement biaisés.
\end{itemize}

\section{Analyse de la Convergence}

\subsection{Accuracy et Loss par Round}
L'évolution de la précision (Accuracy) et de l'erreur (Loss) sur 15 rounds permet de tirer les observations suivantes :
\begin{itemize}
    \item \textbf{Cas IID} : FedAvg, FedAdam et FedProx convergent rapidement (dès le 5ème round) vers une précision d'environ 87\%. L'optimiseur Local SGD stagne aux alentours de 80\%, prouvant l'avantage immédiat de la collaboration.
    \item \textbf{Cas Non-IID} : L'impact de l'apprentissage fédéré est encore plus prononcé. Local SGD souffre fortement de l'hétérogénéité, les modèles locaux sur-apprenant sur des classes majoritaires. FedAvg et FedProx atteignent environ 91\% de précision globale, démontrant une très forte capacité de généralisation malgré les biais locaux.
\end{itemize}

\subsection{Oscillations}
\begin{itemize}
    \item \textbf{FedAvg} : Présente de légères oscillations dans les premiers rounds en environnement non-IID en raison de la divergence des poids locaux.
    \item \textbf{FedAdam} : Réduit les oscillations grâce à son terme de momentum côté serveur, permettant une convergence plus lisse de la fonction de perte (Loss).
    \item \textbf{FedProx} : Limite efficacement la dérive des modèles locaux grâce à son terme proximal ($\mu=0.1$). Cela stabilise l'entraînement, en particulier sur les distributions de données non-IID fortement asymétriques.
\end{itemize}

\section{Communication et Coûts}

\subsection{Nombre de Rounds}
Avec 5 époques locales (epochs) par round, le réseau converge en environ 8 à 10 rounds. Une augmentation du nombre d'époques locales (ex: 10) réduit le nombre de rounds de communication nécessaires pour converger, mais augmente le risque de dérive des clients (client drift) en environnement non-IID, ce qui justifie l'utilisation de FedProx.

\subsection{Taille des Échanges}
Le modèle MLP utilisé est extrêmement léger avec environ 2\,600 paramètres.
En utilisant des flottants 32 bits (4 octets), un échange bidirectionnel (upload + download) coûte environ 21 Ko par hôpital et par round. Pour 15 rounds et 8 hôpitaux, la bande passante totale consommée est inférieure à 3 Mo. Ceci est très efficace et parfaitement adapté aux environnements médicaux contraints en réseau.

\subsection{Ratio Performance / Coût}
L'approche fédérée offre un gain de précision de 7 à 10 points de pourcentage par rapport au modèle Local SGD. Le très faible coût de communication est largement justifié par ce saut de performance et, surtout, par la préservation totale de la confidentialité des données des patients.

\section{Généralisation et Robustesse}

\subsection{Gap Train / Test}
Le sur-apprentissage (overfitting) est naturellement atténué par le processus d'agrégation globale. Bien que les modèles locaux aient tendance à sur-apprendre individuellement (Gap important), le modèle global agrégé conserve une excellente performance sur l'ensemble de test distribué. FedProx montre la meilleure robustesse face au sur-apprentissage local.

\subsection{Robustesse face à la distribution Non-IID}
Les expériences confirment que l'algorithme FedAvg standard peut être ponctuellement affecté par des données fortement asymétriques. En revanche, FedProx et FedAdam se positionnent comme des solutions extrêmement robustes. Ils assurent qu'aucun hôpital (même ceux traitant une écrasante majorité de patients sains ou malades) ne soit laissé pour compte en termes de précision diagnostique.

\section{Conclusion}
L'infrastructure développée et les expériences menées prouvent que l'Apprentissage Fédéré est une solution viable, sécurisée et performante pour le domaine de la santé. Les algorithmes avancés tels que FedProx et FedAdam surpassent la moyenne simple (FedAvg) lorsque les données hospitalières diffèrent fortement. Ils offrent de solides garanties de convergence et une meilleure généralisation, le tout sans jamais avoir besoin de centraliser le moindre dossier médical.

\end{document}
